\documentclass[11pt, a4paper]{article}

\usepackage{lmodern}
\usepackage[T1]{fontenc}
\usepackage[utf8]{inputenc}
\usepackage[french]{babel}
\usepackage{amsmath, amssymb, amsthm}
\usepackage{geometry}
\usepackage{hyperref}

\geometry{margin=2.8cm}

\newtheoremstyle{plain_style}{8pt}{8pt}{\normalfont}{0pt}{\bfseries}{.}{.5em}{}
\theoremstyle{plain_style}
\newtheorem{definition}{D\'{e}finition}
\newtheorem{proposition}{Proposition}
\newtheorem{remarque}{Remarque}
\newtheorem{exemple}{Exemple}

\hypersetup{colorlinks=true, linkcolor=blue, urlcolor=blue}

\title{\textbf{Vers la transform\'{e}e de Laplace}\\
\large Une motivation par les vecteurs propres et les produits scalaires}
\author{}
\date{}

\begin{document}
\maketitle

\section{L'op\'{e}rateur de d\'{e}rivation et ses vecteurs propres}

Consid\'{e}rons l'espace vectoriel $\mathcal{F}$ des fonctions
$f : [0, +\infty) \to \mathbb{C}$ ind\'{e}finiment d\'{e}rivables.
Sur cet espace, l'op\'{e}rateur de d\'{e}rivation
\[
  \frac{d}{dt} : \mathcal{F} \longrightarrow \mathcal{F},
  \qquad f \longmapsto f'
\]
est un op\'{e}rateur lin\'{e}aire. Comme en alg\`{e}bre lin\'{e}aire de dimension finie,
on peut se demander s'il admet des \emph{vecteurs propres}, c'est-\`{a}-dire des fonctions
$\varphi$ telles que
\[
  \frac{d}{dt}\varphi = s\,\varphi
\]
pour un certain scalaire $s \in \mathbb{C}$.

Cette \'{e}quation diff\'{e}rentielle \'{e}l\'{e}mentaire admet pour solution la famille
\[
  \varphi_s : t \longmapsto e^{st}, \qquad s \in \mathbb{C}.
\]
Les exponentielles complexes sont donc exactement les \emph{vecteurs propres} de l'op\'{e}rateur
$d/dt$, associ\'{e}s \`{a} la valeur propre $s$.

\begin{remarque}
En alg\`{e}bre lin\'{e}aire de dimension finie, si un op\'{e}rateur est diagonalisable,
tout vecteur de l'espace se d\'{e}compose en somme de vecteurs propres.
L'ambition de Laplace est analogue~: d\'{e}composer une fonction $f$ arbitraire en une
superposition (continue) d'exponentielles.
\end{remarque}

\section{Sonder $f$ : l'idée du produit scalaire}
Pour d\'{e}tecter la pr\'{e}sence d'une exponentielle $e^{st}$ dans $f$, on s'inspire du
produit scalaire usuel. Dans $\mathbb{R}^n$, le produit scalaire $\langle u, v \rangle$ mesure
l'\emph{alignement} de deux vecteurs~: il est nul lorsqu'ils sont orthogonaux, et maximal
lorsque l'un est multiple de l'autre.

\begin{exemple}
Dans $\mathbb{R}^2$, les vecteurs $(1,0)$ et $(0,1)$ ont un produit scalaire nul~:
ils n'ont rien en commun. En revanche,
$\langle (1,0),\,(2,0) \rangle = 2$~: ils partagent la m\^{e}me direction.
\end{exemple}

Le produit scalaire naturel sur $\mathcal{F}$ est l'int\'{e}grale~:
\[
  \langle f, g \rangle = \int_0^{+\infty} f(t)\,\overline{g(t)}\,dt.
\]
Pour sonder $f$ avec l'exponentielle $e^{st}$, on calcule~:
\[
  \langle f,\, e^{s\,\cdot} \rangle
  = \int_0^{+\infty} f(t)\,e^{\bar{s}\,t}\,dt.
\]
En adoptant la convention usuelle $s \mapsto -s$ (afin de garantir la convergence
pour les parties r\'{e}elles positives de $s$), on obtient le
\emph{prototype de la transform\'{e}e de Laplace}~:
\[
  \boxed{F(s) \;=\; \int_0^{+\infty} f(t)\,e^{-st}\,dt.}
\]

\section{Test du d\'{e}tecteur~: \texorpdfstring{$f(t) = e^{mt}$}{f(t) = e^{mt}}}

Testons notre sonde sur la fonction la plus simple~: $f(t) = e^{mt}$,
avec $m \in \mathbb{C}$. On calcule~:
\[
  F(s) = \int_0^{+\infty} e^{mt}\,e^{-st}\,dt
       = \int_0^{+\infty} e^{(m-s)t}\,dt.
\]
Posons $\alpha = m - s$ et notons $\sigma = \mathrm{Re}(s)$, $\mu = \mathrm{Re}(m)$.
Trois cas se pr\'{e}sentent~:
\begin{align*}
  \sigma > \mu \;(\alpha < 0)
    &:\quad e^{\alpha t} \to 0,\quad
     F(s) = \frac{1}{s-m} \quad\text{(int\'{e}grale convergente),}\\[4pt]
  \sigma < \mu \;(\alpha > 0)
    &:\quad e^{\alpha t} \to +\infty,\quad\text{int\'{e}grale divergente,}\\[4pt]
  s = m \;(\alpha = 0)
    &:\quad e^{\alpha t} = 1,\quad
     \int_0^{+\infty} 1\,dt = +\infty.
\end{align*}

Le troisi\`{e}me cas est le plus r\'{e}v\'{e}lateur. Lorsque $s = m$, l'int\'{e}grande
vaut identiquement $1$ et l'int\'{e}grale diverge~: notre d\'{e}tecteur
\emph{s'affole} exactement au point $s = m$.
La sonde a reconnu l'exponentielle cach\'{e}e dans $f$ et le signale par une divergence.

\begin{proposition}
Lorsque $f(t) = e^{mt}$, la transform\'{e}e de Laplace admet un \emph{p\^{o}le} en $s = m$~:
\[
  F(s) = \frac{1}{s - m}, \qquad \mathrm{Re}(s) > \mathrm{Re}(m).
\]
La position du p\^{o}le encode exactement le param\`{e}tre $m$ de l'exponentielle composant $f$.
\end{proposition}

\begin{remarque}
La condition $\mathrm{Re}(s) > \mathrm{Re}(m)$ d\'{e}finit le \emph{demi-plan de convergence}
de la transform\'{e}e~: la r\'{e}gion du plan complexe dans laquelle la sonde est assez
d\'{e}croissante pour compenser la croissance de $f$.
\end{remarque}

\section{Pourquoi int\'{e}grer sur \texorpdfstring{$[0,+\infty)$}{[0, +infini)} et non sur \texorpdfstring{$\mathbb{R}$}{R}~?}
\paragraph{Causalit\'{e} physique.}
En physique et en ing\'{e}nierie, $t$ repr\'{e}sente le temps~: un syst\`{e}me ne r\'{e}agit
qu'apr\`{e}s avoir re\c{c}u une excitation \`{a} $t = 0$. Int\'{e}grer sur $(-\infty, 0)$
supposerait une m\'{e}moire infinie du pass\'{e}.

\paragraph{Convergence.}
Si $\mathrm{Re}(m) > 0$, alors $f(t) \to +\infty$ quand $t \to -\infty$~:
l'int\'{e}grale $\int_{-\infty}^{0} e^{mt} e^{-st}\,dt$ diverge quelle que soit la valeur de $s$.

\paragraph{Conditions initiales.}
L'int\'{e}gration \`{a} partir de $0$ permet d'incorporer naturellement les conditions
initiales. C'est gr\^{a}ce \`{a} ce choix que l'on obtient la formule cl\'{e}~:
\[
  \mathcal{L}\{f'\}(s) = s\,F(s) - f(0),
\]
qui transforme une \'{e}quation diff\'{e}rentielle en une \'{e}quation alg\'{e}brique en $F(s)$.

\begin{remarque}
La transform\'{e}e de Fourier bil\'{e}at\'{e}rale int\`{e}gre sur $\mathbb{R}$ entier, mais
exige que $f$ soit int\'{e}grable sur $\mathbb{R}$, ce qui exclut les signaux croissants et
toute injection de conditions initiales. Laplace g\'{e}n\'{e}ralise Fourier en autorisant une
croissance exponentielle mod\'{e}r\'{e}e.
\end{remarque}

\section{Bilan et prolongements naturels}

Le raisonnement suivi se r\'{e}sume ainsi~:
\begin{enumerate}
  \item Les exponentielles $t \mapsto e^{st}$ sont les vecteurs propres de $d/dt$.
  \item Pour d\'{e}tecter la composante en $e^{st}$ de $f$, on forme le produit scalaire
        $\displaystyle\int_0^{+\infty} f(t)\,e^{-st}\,dt$.
  \item Sur $f(t) = e^{mt}$, ce produit scalaire diverge exactement en $s = m$,
        signalant la pr\'{e}sence de l'exponentielle.
  \item $F(s)$ encode l'\emph{ADN exponentiel} de $f$~: ses p\^{o}les r\'{e}v\`{e}lent
        les fr\'{e}quences propres.
\end{enumerate}

Les prolongements naturels sont les suivants.

\paragraph{La transform\'{e}e inverse (formule de Bromwich).}
Si $F(s)$ encode $f$, comment reconstruire $f$ depuis $F$~?
\[
  f(t) = \frac{1}{2\pi j}
         \int_{\sigma - j\infty}^{\sigma + j\infty} F(s)\,e^{st}\,ds.
\]
Cette formule boucle la boucle~: elle reconstitue $f$ comme superposition continue
d'exponentielles, exactement l'ambition annonc\'{e}e au d\'{e}part.

\paragraph{Fractions partielles et lecture directe.}
Lorsque $F(s)$ est une fraction rationnelle, la d\'{e}composer en \'{e}l\'{e}ments simples
revient \`{a} lire les exponentielles composant $f$ \`{a} l'\oe{}il nu.

\paragraph{P\^{o}les et stabilit\'{e}.}
La position des p\^{o}les dans le plan complexe r\'{e}v\`{e}le la stabilit\'{e} du syst\`{e}me~:
p\^{o}les \`{a} partie r\'{e}elle n\'{e}gative (syst\`{e}me amorti), sur l'axe imaginaire
(oscillations permanentes), ou \`{a} partie r\'{e}elle positive (divergence).

\section{Injectivité, reconstruction et propriétés opérationnelles}

Maintenant que l'on sait comment la transformée $\mathcal{L}\{f\} = F$ permet de détecter l'ADN de $f$ sous forme d'exponentielles, la question naturelle est de savoir si elle capte suffisamment d'informations pour que l'on puisse reconstruire $f$ connaissant sa transformée de Laplace. La réponse est affirmative et repose sur l'injectivité de l'opérateur.

\subsection{Unicité et reconstruction}

Pour reconstruire $f$ de manière unique, il faut s'assurer que deux fonctions distinctes ne possèdent pas la même transformée. On se place ici dans le cadre des fonctions continues par morceaux et à croissance exponentielle.

\begin{proposition}[Théorème d'unicité de Lerch]
Soient $f$ et $g$ deux fonctions continues sur $[0, +\infty)$ possédant le même demi-plan de convergence. Si, pour tout $s$ dans ce demi-plan, on a $F(s) = G(s)$, alors :
\[
  f(t) = g(t) \quad \forall t \ge 0.
\]
\end{proposition}

\begin{proof}
Par linéarité de l'intégrale, il suffit de montrer que si $F(s) = 0$ pour tout $s > \sigma_0$, alors $f$ est identiquement nulle. Supposons pour simplifier $f$ continue. Quitte à étudier $t \mapsto f(t)e^{-\sigma_1 t}$ avec $\sigma_1 > \sigma_0$, on peut supposer que l'intégrale converge pour $s=0$. On a alors pour tout entier $n \ge 1$ :
\[
  F(n) = \int_0^{+\infty} f(t)e^{-nt}\,dt = 0.
\]
Effectuons le changement de variable $x = e^{-t}$ (donc $dt = -\frac{dx}{x}$). L'intervalle $[0, +\infty)$ devient $(0, 1]$, et l'intégrale se réécrit :
\[
  \int_0^1 f(-\ln x) x^{n-1}\,dx = 0 \quad \forall n \ge 1.
\]
En posant $h(x) = f(-\ln x)$ (prolongée par continuité en $0$), la fonction $h$ est continue sur $[0,1]$ et orthogonale à tous les monômes $x^k$. D'après le théorème d'approximation de Weierstrass, les polynômes sont denses dans l'espace des fonctions continues sur $[0,1]$ pour la norme de la convergence uniforme. Ainsi, $\int_0^1 h(x)^2\,dx = 0$, ce qui implique $h = 0$, et par conséquent $f = 0$.
\end{proof}

\begin{remarque}
L'existence de cette bijection permet d'affirmer que $F(s)$ contient toute l'information de $f(t)$. La formule de reconstruction explicite (inversion de Bromwich) vue précédemment découle de la formule d'inversion de Fourier appliquée à la fonction $t \mapsto f(t)e^{-\sigma t}$.
\end{remarque}

\subsection{Comportement vis-à-vis des opérateurs}

Une seconde question émerge : puisque les exponentielles sont les vecteurs propres de l'opérateur de dérivation $d/dt$, comment la transformée de Laplace se comporte-t-elle vis-à-vis des opérateurs usuels (dérivation, intégration, décalage) ?

\begin{proposition}[Théorème de dérivation et de translation]
Soit $f$ une fonction continûment dérivable et de croissance exponentielle. On a :
\begin{align*}
  \mathcal{L}\{f'\}(s) &= s\,F(s) - f(0) \quad \text{\small (Opérateur de dérivation)} \\
  \mathcal{L}\{t \mapsto f(t)e^{at}\}(s) &= F(s-a) \quad \text{\small (Opérateur de translation fréquentielle)}
\end{align*}
\end{proposition}

\begin{proof}
Pour la dérivation, on utilise une intégration par parties sur l'intervalle $[0, X]$ :
\[
  \int_0^X f'(t)e^{-st}\,dt = \left[ f(t)e^{-st} \right]_0^X - \int_0^X f(t)(-s)e^{-st}\,dt
\]
En développant les bornes :
\[
  \int_0^X f'(t)e^{-st}\,dt = f(X)e^{-sX} - f(0) + s\int_0^X f(t)e^{-st}\,dt
\]
Puisque $s$ est choisi dans le demi-plan de convergence, la fonction $f$ grandit moins vite que $e^{st}$, donc $\lim_{X \to +\infty} f(X)e^{-sX} = 0$. En faisant tendre $X$ vers $+\infty$, on obtient bien :
\[
  \mathcal{L}\{f'\}(s) = s\,F(s) - f(0).
\]

Pour la translation fréquentielle (qui correspond à un produit par une exponentielle dans le domaine temporel), le calcul est direct par substitution dans la définition :
\[
  \mathcal{L}\{f(t)e^{at}\}(s) = \int_0^{+\infty} \left(f(t)e^{at}\right)e^{-st}\,dt = \int_0^{+\infty} f(t)e^{-(s-a)t}\,dt = F(s-a).
\]
\end{proof}

\begin{remarque}
La formule de dérivation montre la puissance de l'approche : l'opérateur de dérivation $d/dt$ (analytique et complexe) est transformé en une simple opération algébrique de multiplication par la variable $s$. C'est exactement le comportement attendu d'un passage dans une base de vecteurs propres.
\end{remarque}

\end{document}